% Faithful typesetting transcription of formal Section4, pp23:24--23:30.
\section*{4 Sharp bootstrapping results from hash functions}
We now show the extremely sharp bootstrapping results for small linear size circuits based on the almost universal hash function constructed above, using a refined kernelization method of Chen, Jin, and Williams [9,10].
We will first prove the general hardness magnification theorem for all sparse NP languages in Section4.1. In Section4.2, we will present stronger hardness magnification results for MCSP, which will utilize the explicitness of our hash construction. Then we will show in Section4.3 that similar techniques can be applied to obtain a bootstrapping result for explicit obstructions (see, e.g., [10]), which formalizes the explicit proofs of circuit lower bounds. In Section4.4, we construct explicit obstructions and prove circuit lower bounds that tightly match these bootstrapping results.
\subsection*{4.1 Hardness magnification for all sparse NP languages}
We first prove the most general version of the hardness magnification result. More discussions are presented after the proof.
\begin{authorclaim}{Theorem4.1}
Let $s=s(n)$ and $T=T(n)$ be two functions such that
\begin{itemize}
\item $\omega(\log n)\le s(n)\le O(\log^2n/\log\log n)$, $s(\Theta(n/\log^2n))=\Theta(s(n))$, $s$ is non-decreasing;
\item $n^\gamma\le T(n)\le2^{O(n)}$, where $\gamma>1$ is an absolute large constant.
\end{itemize}
Then, if there is a $2^{s(n)}$-sparse language $L$ in $\mathsf{NTIME}[T(n)]$, such that $L$ cannot be computed by probabilistic circuits of size $2n+O(ns/\log^2n)$ within error $\exp(-\Omega(s))$, it then follows that $\mathsf{NTIME}[T(2^{O(n^{1/5})})]\not\subseteq\mathsf{SIZE}[2^{cn^{1/5}}]$ for all $c>1$.
\end{authorclaim}
\begin{proof}
Towards a contradiction we assume that $\mathsf{NTIME}[T(2^{O(n^{1/5})})]\subseteq\mathsf{SIZE}[2^{cn^{1/5}}]$ for some $c>1$. Let $L$ be a $2^{s(n)}$-sparse language in $\mathsf{NTIME}[T(n)]$. We now show that $L$ can be computed by probabilistic circuits of size $2n+O(ns/\log^2n)$ within error $\exp(-\Omega(s))$.
Let $H=H^{\mathrm{final}}_{\beta,\ell}=\{H_n\}_{n\ge1}$ be the linear $O(\ell\log^{1+\beta}n)$-samplable $\Theta(\ell)$-output $\exp(-\Omega(\ell))$-almost universal hash function in Theorem3.9 with $\ell(n)=O(s)$ and $\beta=2$, such that the collision probability is at most $2^{-2s}$.\footnote{In particular, assume that $\alpha$ is a constant such that $H$ is $\exp(-\alpha\ell)$-almost universal hash function, then we take $\ell(n)=\lceil2s/\alpha\rceil$.} The circuit complexity of each $h\in\operatorname{supp}(H_n)$ is bounded by $2n+O(ns/\log^2n)$. Since $H$ is efficiently samplable, there is a polynomial time algorithm $M(1^n,v,\cdot)$ computing the hash function corresponding to the seed $v$, where $|v|=O(\ell\log^3n)=O(s\log^3n)$.\footnote{That is, for every $x\in\{0,1\}^n$ and every seed $v$, we have $M(1^n,v,x)=h_v(x)$, where $h_v$ is the hash function corresponding to the seed $v$.}
We now define an intermediate language $L'=\{(n,v,h)\mid\exists x\in\{0,1\}^n\cap L,M(1^n,v,x)=h\}$. For any sufficiently large $n$, we pad the language to make the corresponding length of the tuple in $L'$ has length exactly $m=m(n)=\lfloor(\log(n)/(2c))^5\rfloor$ to obtain $L''$. That is, for sufficiently large $n$ and $v,h$ of appropriate length, we define
\[z_{n,v,h}\triangleq(n,v,h)\Vert1\Vert0^{m(n)-|(n,v,h)|-1}\]
and we have
\[z_{n,v,h}\in L''\Longleftrightarrow(n,v,h)\in L'.\]
Using the straightforward non-deterministic algorithm for $L'$ (guessing $x\in\{0,1\}^n$ and determine whether $M(1^n,v,x)=h$), we can show that $L''\in\mathsf{NTIME}[T(2^{O(m^{1/5})})]\subseteq\mathsf{SIZE}[2^{cm^{1/5}}]$, the second containment follows from our assumption.
The key in our argument is an algorithm for the sparse language $L$ with an oracle access to this intermediate problem $L''$, and then replace the oracle with the small size circuit by assumption. This is formalized in the following lemma.
\begin{authorclaim}{Lemma4.2}
There is a uniform family $D=\{D_n\}_{n\ge1}$ of probabilistic oracle circuits of size $2n+O(ns/\log^2n)$, such that for every input length $n$, the followings hold:
\begin{enumerate}
\item Every $D\in\operatorname{supp}(D_n)$ contains at most one $L''$ oracle gate of fan-in $\lfloor\log^5(n)/(32c^5)\rfloor$.
\item $D$ decides the language $L$ with error at most $2^{-s}$.
\end{enumerate}
\end{authorclaim}
\begin{proof}
Consider the following probabilistic circuit family $D=\{D_n\}_{n\ge1}$. Given input $x$ of length $n$, $D_n$ is constructed as follows: we randomly choose a seed $v$ to sample a hash function $h_v$ from $H_n$; and then query whether $z_{n,v,h_v(x)}$ is in $L''$ via the oracle gate, where $h_v$ is the hash function corresponding to the seed $v$. The circuit complexity of $D_n$ equals the circuit complexity of the hash function $h_v$, which is at most $2n+O(ns/\log^2n)$ by Theorem3.9, and every $D\in\operatorname{supp}(D_n)$ only needs to call the $L''$ oracle once with input length $m=\lfloor\log^5(n)/(32c^5)\rfloor$. Clearly, by the definition of $L''$, for any $x\in L$, $D_n$ accepts $x$ with probability 1. For any $x\notin L$, we have
\begin{align*}
\Pr_v[D_n\text{ accepts }x]
&=\Pr_{h\leftarrow H_n}[\exists x'\in L\cap\{0,1\}^n\text{ s.t. }h(x)=h(x')]\\
&\le\sum_{x'\in L\cap\{0,1\}^n}\Pr_{h\leftarrow H_n}[h(x)=h(x')]\\
&=2^s\cdot2^{-2s}\le2^{-s}.
\end{align*}
\end{proof}
Now applying Lemma4.2 and replacing the oracle query by circuits using the fact that $L''\in\mathsf{SIZE}[2^{cn^{1/5}}]$, we finish the proof of Theorem4.1.
\end{proof}
We now discuss some typical choices of parameters in the above theorem. Let the sparsity parameter $s(n)=\log^2n/\log\log n$.\footnote{For a typical $2^{s(n)}$-sparse language in NP, consider $\mathrm{MCSP}[n^{1.9}]$ on input length $N=2^n$.}
\begin{itemize}
\item A standard form of hardness magnification result similar to the one in [9] but quantitatively stronger can be obtained by choosing $T(n)=\operatorname{poly}(n)$. Then the theorem says that, if there exists an $n^{\log n/\log\log n}$-sparse language in NP that does not have probabilistic circuits of size $2n+O(n/\log\log n)$, then $\mathsf{NTIME}[2^{O(n^{1/5})}]\not\subseteq\mathsf{SIZE}[2^{cn^{1/5}}]$ for all $c>0$. By a padding argument, this implies that $\mathsf{NP}\not\subseteq\mathsf{SIZE}[n^c]$ for all $c>0$. This establishes Theorem1.2.
\item Being less ambitious, we can let $T(n)=2^{n^{o(1)}}$ and show that, even the existence of an $n^{\log n/\log\log n}$-sparse language in $\mathsf{NTIME}[2^{n^{o(1)}}]$ that does not have probabilistic circuits of size $2n+O(n/\log\log n)$, would be enough to imply that $\mathsf{NTIME}[2^{2^{o(n^{1/5})}}]\not\subseteq\mathsf{SIZE}[2^{cn^{1/5}}]$, and hence $\mathsf{NTIME}[2^{n^{o(1)}}]\not\subseteq\mathsf{SIZE}[n^c]$ for all $c>0$ and $\mathsf{NEXP}\not\subseteq\mathsf P/\mathrm{poly}$, which would already be a major breakthrough in circuit complexity. This establishes Theorem1.3.
\end{itemize}
\subsection*{4.2 Hardness magnification for MCSP}
We now utilized the strongly explicit hash function constructed above to obtain a stronger hardness magnification theorem for MCSP. Intuitively, this is possible since the yes-instances of MCSP can be efficiently encoded in a much shorter string. Indeed, for $\mathrm{MCSP}[s(n)]$, one only need $O(s(n)\log s(n))$ bits to encode a small circuit, instead of $2^n$ bits for a whole truth table.
We assume a paddable encoding of circuits, i.e., any string $x$ and $x\Vert0$ encode the same circuit. This can be done with only a constant overhead. For a circuit $C$, we use the notation $\langle C\rangle$ to denote its encoding.
\begin{authorclaim}{Theorem4.3}
Let $n\le s(n)\le O(n^2/\log^2n)$ be a non-decreasing size parameter that satisfies $s(\log(\Theta(n/\log^2n)))=\Theta(s(\log n))$. Let $N=2^n$ be the truth table length. Let $g=g(n)=s(n)\log s(n)$. If $\mathrm{MCSP}[s(n)]$ cannot be computed by probabilistic circuits of size $2N+O(Ng/\log^2N)$ within error $\exp(-\Omega(g))$, then there is some $c\in(0,1)$ such that $\oplus\mathsf P\not\subseteq\mathsf{SIZE}[2^{N^c}]$.
\end{authorclaim}
\paragraph{Intuition.} The proof is similar to the one for Theorem4.1, but we need to make the oracle from Lemma3.1 computable in $\oplus\mathsf P$ instead of in super-polynomial non-deterministic time. Concretely, we need to slightly modify the intermediate problem so that we can make use of the strongly explicitness of our hash construction $H^{\mathrm{final}}$.
\begin{proof}
Towards a contradiction, suppose that $\oplus\mathsf P\subseteq\mathsf{SIZE}[2^{N^c}]$ for all $c\in(0,1)$, we now prove that $\mathrm{MCSP}[s(n)]$ can be computed by probabilistic circuits of size $2N+O(Ng/\log^2N)$ within error $2^{-\Omega(g(n))}$.
Let $N=2^n$ be the truth table length of $n$-input functions. Take $H=H^{\mathrm{final}}_{\beta,\ell}=\{H_N\}_{N\ge1}$ to be the linear $O(\ell\log^{1+\beta}N)$-samplable $\Theta(\ell)$-output $\exp(-\Omega(\ell))$-almost universal hash function in Theorem3.9 with $\ell(N)=O(g(n))=O(g(\log N))$ and $\beta=2$, such that the collision probability is at most $2^{-2g}$.\footnote{Note that since $s(n)\le O(n^2/\log^2n)$, we have $\ell(N)=O(\log^2N/\log\log N)$ and $\ell$ is non-decreasing. Hence, $\ell$ satisfies the requirements in Theorem3.9.}
We denote the hash function of input length $N$ corresponding to the seed $v$ as $H_{N,v}(\cdot)$. By the strongly explicitness (Corollary3.15), there exists a algorithm $M(N,v,i,j)$ running in $\operatorname{polylog}(N)=\operatorname{poly}(n)$ time that decides whether the $i$th output bit depends on the $j$th input bit in $H_{N,v}$. Note that for a circuit $C$ of input length $n$,\footnote{Here we use the notation $\mathrm{tt}(C)$ to represent the truth table of $C$. Particularly, if $C$ is a single-output circuit of $n$ inputs, then $\mathrm{tt}(C)$ is a string of length $2^n$ such that $\mathrm{tt}(C)_i=C(i)$.}
\[H_{N,v}(\mathrm{tt}(C))_i=\sum_{j=0}^{N-1}M(N,v,i,j)\cdot C(j)\pmod2,\]
and $M(N,v,i,j)\cdot C(j)$ can be computed in $\operatorname{poly}(n,|\langle C\rangle|)$ time, so the language $L^h=\{(N,i,v,\langle C\rangle)\mid H_{N,v}(\mathrm{tt}(C))_i=1\}$ is decidable in $\oplus\mathsf P$. Hence if we define an intermediate language $L'=\{(N,v,h)\mid\exists C\in\mathsf{SIZE}[s(n)],H_{N,v}(\mathrm{tt}(C))=h\}$ similar to the proof of Theorem4.1, it is now decidable in NP given oracle access to $\oplus\mathsf P$.
Let $m=O(\ell\log^3n)$ be the input length of $L'$. Note that $L'\in\mathsf{NP}^{\oplus\mathsf P}\subseteq\mathsf{BPP}^{\oplus\mathsf P}\subseteq\mathsf P_{/\mathrm{poly}}^{\oplus\mathsf P}\subseteq\oplus\mathsf P/\mathrm{poly}$, where $\mathsf{NP}^{\oplus\mathsf P}\subseteq\mathsf{BPP}^{\oplus\mathsf P}$ follows from [42] (see also [15]) and $\mathsf P_{/\mathrm{poly}}^{\oplus\mathsf P}\subseteq\oplus\mathsf P/\mathrm{poly}$ follows from $\oplus\mathsf P^{\oplus\mathsf P}\subseteq\oplus\mathsf P$ [37].\footnote{Since $\oplus\mathsf P^{\oplus\mathsf P}\subseteq\oplus\mathsf P$ implies that $\mathsf P^{\oplus\mathsf P}\subseteq\oplus\mathsf P$, it follows that the evaluation of polynomial-size circuits with $\oplus\mathsf P$ oracles can be computable in $\oplus\mathsf P$, which further implies that $\mathsf P_{/\mathrm{poly}}^{\oplus\mathsf P}\subseteq\oplus\mathsf P/\mathrm{poly}$.} Since $\oplus\mathsf P\subseteq[2^{N^c}]$ for all $c\in(0,1)$ by the assumption, we know that the evaluation of $\oplus\mathsf P/\mathrm{poly}$ circuits can be computable in $\mathsf{SIZE}[2^{N^c}]$ for all $c\in(0,1)$, which implies that $L'\in\mathsf{SIZE}[\sqrt N]$. The rest of the proof follows similar to Theorem4.1 and the fact that the $m=O(\ell\log^3N)=o(\log^5N)$.
\end{proof}
\subsection*{4.3 Explicit obstruction}
We now formally state and prove our results regarding explicit obstructions. We begin by formally defining the notion of explicit obstructions.
\begin{authorclaim}{Definition4.4 (Explicit obstruction)}
An explicit obstruction of size $S(n)$ computable in $C$ against $D$ is a family of lists of input-output pairs $O=\{O_n\}_{n\ge1}$ satisfying the following.
\begin{itemize}
\item $|O_n|\le S(n)$ for all sufficiently large $n$.
\item There is a machine in $C$ that prints the set $O_n$ given input $1^n$.
\item For every $n$, $O_n=\{(x_i,y_i)\}$ satisfies that $x_i\ne x_j$ for all $i\ne j$.
\item For all sufficiently large $n$, and for every $n$-input $D$ circuit $f$, there is a pair $(x_i,y_i)\in O_n$ such that $f(x_i)\ne y_i$.
\end{itemize}
\end{authorclaim}
\begin{authorclaim}{Theorem4.5}
There is an absolute constant $\gamma>0$ such that the following holds. Let $n^\gamma\le T(n)\le2^n$ and $\log n\le s(n)\le\min\{O(\log^2n/\log\log n),\log T(n)\}$ being non-decreasing and satisfying $s(\Theta(n/\log^2n))=\Theta(s(n))$. If there is an explicit obstruction of size $2^{s(n)}$ computable in $\mathsf{DTIME}[T(n)]$ against $2n+O(n/\log\log n)$-size circuits, then $\mathsf{DTIME}[T(2^{O(n^{1/5})})]\not\subseteq\mathsf{SIZE}[2^{cn^{1/5}}]$ for all $c>1$.
\end{authorclaim}
\begin{proof}
Let $t(n)=2^{s(n)}$. Towards a contradiction we assume that there is some constant $c>1$, such that $\mathsf{DTIME}[T(2^{O(n^{1/5})})]\subseteq\mathsf{SIZE}[2^{cn^{1/5}}]$. Suppose that $O=\{O_n\}_{n\ge1}$ is an explicit obstruction against $2n+O(n/\log\log n)$ size circuits. For any sufficiently large input length $n$, we suppose that $O_n=\{(x_{n,1},y_{n,1}),(x_{n,2},y_{n,2}),\ldots,(x_{n,t(n)},y_{n,t(n)})\}$. Our goal is to design a circuit with extremely small size, but agree with $O_n$ on all its input-output pairs. Following the similar proof outline as for Theorem4.1, we begin by using an almost universal hash function to kernalize the inputs from $O_n$.
Let $H=H^{\mathrm{final}}_{\ell,\beta}=\{H_n\}_{n\ge1}$ be the linear $O(\ell\log^{1+\beta}n)$-samplable $\Theta(\ell)$-output $\exp(-\Omega(\ell))$-almost universal hash function in Theorem3.9 with $\ell(n)=O(s)$ and $\beta=2$ such that the collision probability is at most $2^{-3s}$. Let $n$ be a sufficiently large input length. We call a hash function $h$ good if it is perfect on the inputs in $O_n$, i.e., any two distinct inputs in $O_n$ have different hash values assigned by $h$. For a randomly chosen function $h\leftarrow H_n$, the probability of it not being good is bounded by
\begin{align*}
\Pr_{h\leftarrow H_n}[\exists1\le i<j\le t(n)\text{ s.t. }h(x_{n,i})=h(x_{n,j})]
&\le\sum_{1\le i<j\le t(n)}\Pr_{h\leftarrow H_n}[h(x_{n,i})=h(x_{n,j})]\\
&\le\binom{2^s}2 2^{-3s}\le2^{-s}.
\end{align*}
So a good hash function always exists in the support of $H_n$ for all sufficiently large $n$. For any large input length $n$, we arbitrarily fix such a good hash and denote its seed by $v_n^{\mathrm{good}}$.
Let $h_v$ be the hash function corresponding to the seed $v$. We define an intermediate language $L'=\{(n,v,h)\mid\exists1\le i\le t(n),y_{n,i}=1\wedge h_v(x_{n,i})=h\}$. We again pad its input to have length exactly $m=\lfloor(\log(n)/(2c))^5\rfloor$, and form a padded language $L''$. Then $L''\in\mathsf{DTIME}[T(2^{O(m^{1/5})})]\subseteq\mathsf{SIZE}[2^{cm^{1/5}}]$. We also note that $2^{cm^{1/5}}\le\sqrt n$.
Then we consider the function $f_n(x)\triangleq L'(n,v_n^{\mathrm{good}},H_n(v_n^{\mathrm{good}},x))$, where $H_n(v,\cdot)$ is the hash function in $H_n$ corresponding to the seed $v$. By an argument similar to that of Theorem4.1, we can show that $f_n$ can be decided by a circuit of size $2n+O(n/\log\log n)+\sqrt n=2n+O(n/\log\log n)$, but totally agrees with $O_n$, contradicting to the assumption that $O$ is an explicit obstruction against circuits of such size.
\end{proof}
\subsection*{4.4 Unconditional lower bounds for sparse languages}
Now we complement the results mentioned in previous subsections by constructing an explicit obstruction against $B_2$ circuits of size $2n-O(1)$ and a corresponding sparse language in P with a $2n-O(1)$ probabilistic circuit lower bound. They form sharp bootstrapping thresholds together with Theorem4.1 and Theorem4.5.
The main idea behind our explicit obstruction is the investigation of a combinatorial structure in Boolean circuits called critical path introduced in [13], which was used to prove $2n-O(1)$ circuit lower bounds for PRFs and hash functions.
For the simplicity of presentation, we assume without loss of generality that our circuits are normalized, in the sense that there is no non-output gates with out-degree 0. This is without loss of generality since redundant gates with out-degree 0 can be removed.
\begin{authorclaim}{Definition4.6 (Critical path)}
Let $C$ be a circuit, and $u$ be an input variable of it. The critical path of $u$ in $C$ is a sequence of vertices $v_0,v_1,\ldots,v_k$ satisfying the following conditions:
\begin{enumerate}
\item $v_0=u$, and $v_i$ is a descendent of $v_{i-1}$ for all $i\ge1$, and
\item $\operatorname{out\text{-}degree}(v_i)=1$ for all $0\le i<k$, and $\operatorname{out\text{-}degree}(v_k)\ne1$.
\end{enumerate}
\end{authorclaim}
Fix a circuit with $n$ input bits, we can obtain a total of $n$ critical paths. The crucial observation in [13] is that, if all critical paths do not intersect with one another, then the circuit must contain at least $2n-O(1)$ gates. This is formalized below.
\begin{authorclaim}{Lemma4.7 ([13], Lemma6.4)}
For any normalized $n$-input single-output circuit $C$ with no intersecting critical paths and no input variables with out-degree 0, the number of gates in the circuit is at least $2n-2$.
\end{authorclaim}
Let $O=\{O_n\subseteq\{0,1\}^n\times\{0,1\}\}_{n\ge1}$, where
\[O_n=\{(x,0)\mid|x|\in\{0,2,n-2,n-1\}\}\cup\{(x,1)\mid|x|\in\{1,n\}\}.\]
In the remaining part of the section, we will first prove that $O$ is an explicit obstruction against $2n-2$ size circuits, and then present a general connection between explicit obstruction and probabilistic circuit lower bounds.
\paragraph{Explicit obstruction.} According to Lemma4.7, we only need to prove that any circuit with intersecting critical paths or input variables with out-degree 0 does not fully agree with $O$. For circuits with input variables of out-degree 0, this can be verified straightforwardly. Now we prove the case for circuits with intersecting critical paths.
\begin{authorclaim}{Lemma4.8}
For any circuit $C_n:\{0,1\}^n\to\{0,1\}$ with intersecting critical paths, there must exist a pair $(x,b)\in O_n$ such that $C_n(x)\ne b$.
\end{authorclaim}
\begin{proof}
Suppose that the critical paths of $u$ and $v$ in $C_n$ intersect. Let $G$ be the first gate on the intersection. Let $f_G$ be the function computed by $G$. Assume that we take a restriction $\rho$ to all variables except $u$ and $v$, and consider the restricted function $C_n\upharpoonright\rho$. The key observation is that, for any $\rho$, there are functions $\phi_\rho,\psi_\rho,\chi_\rho:\{0,1\}\to\{0,1\}$ such that
\[C_n\upharpoonright\rho(u,v)=\chi_\rho(f_G(\phi_\rho(u),\psi_\rho(v))).\]
Based on the characterization of different functions in $B_2$, we call a function $f$ quadratic if it has the form $f(u,v)=((u\oplus c_1)\wedge(v\oplus c_2))\oplus c_3$. We call a function linear if it has the form $f(u,v)=u\oplus v\oplus c$. The pivotal point in [13] is that for a fixed pair of $u$ and $v$, $C_n\upharpoonright\rho$ cannot be quadratic under some restriction $\rho_1$, then become linear under another restriction $\rho_2$. However, by checking the truth table it can be verified that our construction of $O_n$ forces $C_n\upharpoonright\rho$ to be the linear function $(u,v)\mapsto u\oplus v$ for all-zero restriction, and the quadratic function $(u,v)\mapsto u\wedge v$ for all-one restriction, which leads to contradiction.
\end{proof}
\begin{authorclaim}{Corollary4.9}
$O$ is an explicit obstruction against $2n-2$ size $B_2$ circuits.
\end{authorclaim}
\paragraph{Probabilistic circuit lower bound.} We now present a general reduction from explicit obstruction to probabilistic circuit lower bounds for sparse languages in P.
\begin{authorclaim}{Lemma4.10}
For any circuit class $C$, if there exists an explicit obstruction $O$ of size $S(n)$ against $C$, then every language agreeing with $O$ cannot be computed by probabilistic $C$-circuits with error probability less than $1/S(n)$, even infinitely often.
\end{authorclaim}
\begin{proof}
Towards contradiction, let $L$ be a language agreeing with $O$, and assume that there exists a probabilistic $C$-circuit $C$ computing $L$ with error probability smaller than $1/S(n)$, for infinitely many $n$. Then for those $n$, we have
\[\mathbb E_C[|\{(x,b)\in O_n\mid C(x)\ne b|\}]<1.\]
By the averaging argument, there exists a deterministic circuit that agrees with all pairs in $O_n$. This guarantees the existence of a family of circuit in $C$ that agrees with $O$ on those $n$, which leads to a contradiction.
\end{proof}
In particular, let $L=\{L_n\}_{n\ge1}$ be a language such that $L_n\triangleq\{x\in\{0,1\}^n\mid(x,1)\in O_n\}$, then it is a sparse language in P that agrees with $O$. Together with Corollary4.9, we obtain the following lower bound.
\begin{authorclaim}{Corollary4.11}
There exists a $O(n)$-sparse language in P that cannot be computed (even infinitely often) by probabilistic $B_2$ circuits of size $2n-2$ with error probability smaller than $1/(2n^2)$.
\end{authorclaim}
\begin{authorclaim}{Remark4.12}
We note that the error probability in the lower bound cannot be trivially boosted to $1/3$, since the complexity overhead of error reduction is not affordable when we are dealing with small linear-size circuits. Therefore it might be significantly more difficult to prove a similar circuit lower bound with constant error probability.
\end{authorclaim}
